\documentclass[11pt]{article}
\usepackage{fontspec}
\usepackage[french,english]{babel}
\usepackage{amsmath}
\usepackage{amssymb}
\usepackage{amsthm}
\usepackage{mathrsfs}
\usepackage{tikz-cd}
\usepackage{enumitem}
\usepackage{lmodern}

\newcommand{\printedpage}[1]{\marginpar{\footnotesize p.~#1}}
\newcommand{\msp}[1]{\marginpar{\footnotesize [#1]}}
\newcommand{\Ob}{\operatorname{Ob}}
\newcommand{\Filt}{\mathrm{Filt}}
\newcommand{\Ind}{\operatorname{Ind}}
\newcommand{\card}{\operatorname{card}}
\newcommand{\cardred}{\operatorname{card\,red}}

\begin{document}

\noindent\printedpage{59}% head: A. Grothendieck : Dérivateurs, Ch. XVIII, Édition des Universités de Montpellier II et Paris VII
% foot: Version 17 décembre 2025
except that one must take care with the cardinal bound, $\card I \leqslant \pi'^{\pi}$. But if $F \in \Ob\mathcal{F}^{\pi'}$, since by (4.7.6.1)
\[
\mathcal{F}^{\pi'} \;\subset\; \underbrace{\mathcal{M}_{\pi'}}_{\subset\mathcal{M}_{\pi'^{\pi}}\subset\Filt'^{\pi'^{\pi}}} \;\subset\; \Filt'^{\pi'^{\pi_0}} ,
\]
hence a fortiori
\[
\mathcal{F}^{\pi'} \;\subset\; \Filt^{\pi'^{\pi}} \underbrace{\;=\;}_{\text{if } \pi'^{\pi}\geqslant\pi_1} \Filt^{\pi'^{\pi}} , % [sic] first \Filt^{\pi'^{\pi}} -> \Filt'^{\pi'^{\pi}}, otherwise the equality under the brace would be trivial (cf. the preceding display)
\]
whence, by 4.7.3, with $\pi'$ replaced by $\pi'^{\pi}=\pi''$, the possibility of taking $I$ with
\[
\card I \;\leqslant\; \pi''^{\pi} \;=\; (\pi'^{\pi})^{\pi} \;=\; \pi'^{\pi\cdot\pi} \;=\; \pi'^{\pi} ,
\]
which completes the proof.

\noindent\msp{page 175}\textbf{Remark 4.7.7.2.} I was obliged, in the statement of condition c), to bring in a second cardinal $\pi_1$, and this for lack of knowing whether, for $\pi''$ such that $\pi''=\pi''^{\pi_0}$, $\Filt^{\pi''}$ is stable under direct factors. (To do things properly, one would have to either prove that it is stable, hence that we have in 4.7.1 c $\Longrightarrow$ b, i.e. $\mathcal{M}_\pi\subset\Filt^{\pi}$ (and not only c $\Longrightarrow$ b$'$, i.e. $\mathcal{M}_\pi\subset\Filt'^{\pi}$) if $\pi=\pi^{\pi_0}$, or else make a counterexample!) This somewhat strange ingredient, we shall alas be obliged to drag it along in the definition of generalized cardinal filtrations, if we do not want to forbid ourselves to regard in the present case $(\mathcal{F}^{\pi})_{\pi\geqslant\pi_0}$ as such a filtration, and to allow ourselves only to regard $(\mathcal{F}^{\pi})_{\pi\geqslant\pi_1}$ as such, which would seem ridiculous to me. In practice (even if theoretically the introduction of $\pi_1$ in the description of a generalized cardinal filtration were indeed necessary) we shall hardly have cases where it is necessary to make mention of $\pi_1$ in condition c. Indeed, $\pi\geqslant\pi_0$ and $\pi^{\pi_0}=\pi$ imply $\pi\geqslant 2^{\pi_0}$, and then $\pi'^{\pi}\geqslant 2^{(2^{\pi_0})}$. If therefore $\pi_1\leqslant 2^{(2^{\pi_0})}$, i.e. $\card C_0\leqslant 2^{(2^{\pi_0})}$, then the condition $\pi^{\pi'}\geqslant\pi_1$ in c) is automatically satisfied, and there is therefore no need to mention it. Now it is very rare to have to deal with small categories $C_0$, whose cardinality is $>$ the power % [sic] faut -> faute, "faute de savoir" intended; [sic] \pi^{\pi'}\geqslant\pi_1 -> \pi'^{\pi}\geqslant\pi_1, the condition is written \pi'^{\pi}\geqslant\pi_1 elsewhere (cf. 4.7.8)
\msp{page 176}
of the continuum $c$, and a fortiori, such that $\card C_0 > 2^{c}$. If on the other hand $\card C_0\leqslant 2^{c}$, since $2^{c}\leqslant 2^{(2^{\pi_0})}$ (since $\pi_0\geqslant\aleph_0$), we shall indeed have $\pi_1\leqslant 2^{(2^{\pi_0})}$, and we can rest easy.

\textbf{Definition 4.7.8.} Let $\mathcal{M}$ be a $\mathfrak{U}$-category. We call a $\mathfrak{U}$-\textit{generalized cardinal filtration} (or simply generalized cardinal filtration, when $\mathfrak{U}$ is understood) the datum of an infinite cardinal $\pi_0$, and of a filtration $(\mathcal{F}^{\pi})_{\pi\geqslant\pi_0}$ of $\mathcal{M}$ by strictly full subcategories $\mathcal{F}^{\pi}$, indexed by the $\mathfrak{U}$-cardinals $\geqslant\pi_0$, and satisfying conditions a) b) c) stated in 4.7.7, where in c) we put
\[
\pi_1 \;=\; \sup(\pi_0,\cardred\mathcal{F}^{\pi_0}) \, . \tag*{(4.7.8.1)}
\]
(Thus if $\cardred\mathcal{F}^{\pi_0}\leqslant 2^{(2^{\pi_0})}$, there is no need to take account of the condition $\pi'^{\pi}\geqslant\pi_1$, cf. 4.7.7.1.) We say that the generalized cardinal filtration (fcg) is \textit{accessible}, if it satisfies in addition condition d). We say that it is \textit{canonical} if we even have
\[
\mathcal{F}^{\pi} \;=\; \mathcal{M}_{\pi} \tag*{(4.7.8.2)}
\]
for every $\pi\geqslant\pi_0$.

\textbf{Proposition 4.7.9.} \textit{Let $(\mathcal{F}^{\pi})_{\pi\geqslant\pi_0}$ be an fcg on the $\mathfrak{U}$-category $\mathcal{M}$, and $\pi>\pi_0$ a cardinal such that $\pi^{\pi_0}=\pi$. Then we have}
\printedpage{60}% head: A. Grothendieck : Dérivateurs, Ch. XVIII, Édition des Universités de Montpellier II et Paris VII
% foot: Version 17 décembre 2025
\msp{page 177}
\[
\left\{
\begin{aligned}
&\mathcal{M}_\pi\subset\mathcal{F}^{\pi}\quad\text{\textit{(if $\pi=\pi^{\pi_0}$),}}\\
&\text{\textit{hence} }\mathcal{M}_\pi=\mathcal{F}^{\pi}\text{ \textit{if moreover the fcg is} accessible.}
\end{aligned}
\right.
\tag*{(4.7.9.1)}
\]

\textsc{Proof} (\footnote[75]{\texttt{[In what follows, $\mathcal{F}^{\pi}$ has sometimes been written $\mathcal{F}_{\pi}$, often corrected to $\mathcal{F}^{\pi}$ by the author. We write $\mathcal{F}^{\pi}$ throughout.]}}). Applying condition c) for fcg to an $F\in\Ob\mathcal{M}_\pi$, we find that $F$ is a direct factor of an $x_i\in\Ob\mathcal{F}^{\pi}$. Since $\mathcal{F}^{\pi}$ is stable under $\varinjlim$ large relative to $\pi_0$ and of cardinality $\leqslant\pi$, and since \textit{in addition} $\pi>\pi_0$, we conclude that:

\textbf{Lemma 4.7.9.1} \texttt{[twice 4.7.9.1]}. \textit{If $\pi>\pi_0$ (\footnote[76]{But $\mathcal{F}^{\pi_0}$ is not necessarily stable under direct factors, as we see by taking $\mathcal{M}=\Ind_{\pi_0}(C)$, $C$ a small non-Karoubian category, and the cardinal filtration for the $\Filt'^{\pi}$ -- we note that $\Filt'^{\pi_0}=\Filt^{\pi_0}\simeq e$, since a category large relative to $\pi_0$ and of cardinality $\leqslant\pi_0$ is cofinal to $e$.}), $\mathcal{F}^{\pi}$ is stable under direct factors.}

Whence it follows that $F\in\mathcal{F}^{\pi}$, whence (4.7.9.1).

\textbf{Corollary 4.7.10.} \textit{Suppose the fcg accessible. For every cardinal $\pi\geqslant\pi_0$, we have}
\[
\mathcal{F}^{\pi} \;\subset\; \mathcal{M}_\pi \;\subset\; \mathcal{F}^{\pi^{\pi_0}} \; ( \;=\; \mathcal{M}_{\pi^{\pi_0}}) \, . \tag*{(4.7.10.1)}
\]

Indeed, $\mathcal{F}^{\pi}\subset\mathcal{M}_\pi$ by the definition of accessible fcg, and $\mathcal{M}_\pi=\mathcal{M}_{\pi^{\pi_0}}$ (since $\pi\leqslant\pi^{\pi_0}$), now $\mathcal{M}_{\pi^{\pi_0}}=\mathcal{F}^{(\pi^{\pi_0})}$ by virtue of (4.7.9.1) applied to $\pi'=\pi^{\pi_0}$ instead of $\pi$ (since $\pi'^{\pi_0}=\pi'$), whence (4.7.10.1). We shall note that this inclusion generalizes (4.7.9.1), it gives $\mathcal{F}^{\pi}=\mathcal{M}_\pi$ if $\pi=\pi^{\pi_0}$.

\textbf{4.7.11.} It seems to me that the only useful generalized cardinal filtrations are those which are \textit{accessible.} In what follows, I am going to confine myself to this case. Since in that case $\mathcal{F}^{\pi}\subset\mathcal{M}_\pi$, we conclude from this
\msp{page 178}
that the canonical functor
\[
\Ind_\pi(\mathcal{F}^{\pi}) \;\hookrightarrow\; \mathcal{M} \tag*{(4.7.11.1)}
\]
extending the inclusion $\mathcal{F}^{\pi}\subset\mathcal{M}$, so as to commute with $\varinjlim$ large relative to $\pi$, is \textit{fully faithful.} Moreover, when $\pi^{\pi_0}=\pi$, condition c) for fcg tells us among other things that this functor is essentially surjective, hence that it is an \textit{equivalence of categories}:
\[
\Ind_\pi(\mathcal{F}^{\pi}) \xrightarrow{\ \sim\ } \mathcal{M} \qquad (\text{if } \pi=\pi^{\pi_0}). \tag*{(4.7.11.2)}
\]

\textbf{Proposition 4.7.12.} \textit{Let $(\mathcal{F}^{\pi})_{\pi\geqslant\pi_0}$ be an accessible fcg (fcga) on the $\mathfrak{U}$-category $\mathcal{M}$, and let $\pi$ be a cardinal $\geqslant\pi_0$. The following conditions are equivalent.}
\begin{enumerate}[label=\textit{\alph*)}]
\item \textit{The fully faithful functor (4.7.11.1) is essentially surjective, i.e. is an equivalence of categories.}
\item \textit{For every $F\in\Ob\mathcal{M}$, $\mathcal{F}^{\pi}/F$ is a filtering category and large relative to $\pi$.}
\end{enumerate}

\noindent\printedpage{61}% head: A. Grothendieck : Dérivateurs, Ch. XVIII, Édition des Universités de Montpellier II et Paris VII
% foot: Version 17 décembre 2025
\textit{When these conditions are satisfied, and if $\pi>\pi_0$, we have}
\[
\mathcal{F}^{\pi} \;=\; \mathcal{M}_\pi \, . \tag*{(4.7.12.1)}
\]

\textsc{Proof}. The equivalence of a) and b) follows at once from 4.4.2 and from the definition of categories large relative to $\pi$ (p.~81). That these
\msp{page 179}
conditions imply $\mathcal{F}^{\pi}=\mathcal{M}_\pi$, i.e. $\mathcal{M}_\pi\subset\mathcal{F}^{\pi}$ in the case where $\pi>\pi_0$, is seen as for (4.7.9.1) by using Lemma 4.7.9.1.

\textbf{Definition 4.7.13.}
\begin{enumerate}[label={$\arabic*^{\circ}$)}]
\item A cardinal $\pi\geqslant\pi_0$ is called a \textit{useful cardinal} for the fcga $(\mathcal{F}^{\pi})_{\pi\geqslant\pi_0}$, if it satisfies the equivalent conditions a) and b) of 4.7.12, and if in addition in the case where $\pi=\pi_0$, we have $\mathcal{F}^{\pi}=\mathcal{M}_\pi$. (Thus if $\pi$ is a useful cardinal, we always have $\mathcal{F}^{\pi}=\mathcal{M}_\pi$.) In the contrary case, we say that $\pi$ is a \textit{superfluous cardinal} (for the given fcga). % [sic] le fcga donnée -> la fcga donnée, gender agreement (feminine)
\item A cardinal $\pi$ is called \textit{regular} for the given fcg, if it is useful, and if in addition the cardinal bound stated in (4.7.7.1) is valid.
\end{enumerate}

\textbf{4.7.14} Thus, all cardinals $\pi\geqslant\pi_0$ such that $\pi^{\pi_0}=\pi$ are useful, and even regular for the given \texttt{[generalized]} cardinal filtration. Let us note that if all cardinals $\pi\geqslant\pi_0$ are useful, then we have $\mathcal{F}^{\pi}=\mathcal{M}_\pi$ for every $\pi\geqslant\pi_0$, hence the fcga is canonical. Let us also note that in all cases, if $(\mathcal{F}^{\pi})_{\pi\geqslant\pi_0}$ is an fcga, then $(\mathcal{M}_\pi)_{\pi\geqslant\pi_0}$ is one too.

\noindent\msp{page 180}Indeed, condition a) for fcg (cf. 4.7.7) follows from the same condition for $(\mathcal{F}^{\pi})$, and from (4.7.10.1), $\mathcal{M}_\pi\subset\mathcal{F}^{\pi^{\pi_0}}$. Condition b) is a general property of $\pi$-accessibility, taking into account that $(\mathcal{F}^{\pi_0})_0$ is a small generating subcategory consisting of $\pi_0$-accessible objects of $\mathcal{M}$. Condition d) is satisfied tautologically, and it remains to check condition c). But we have seen that for $\pi=\pi^{\pi_0}$, we have $\mathcal{M}_\pi=\mathcal{F}^{\pi}$ (4.7.9.1), hence condition c) for $(\mathcal{M}_\pi)$ follows from the same condition for $(\mathcal{F}^{\pi})$. Thus, we find:

\textbf{Proposition 4.7.15.} \textit{If $(\mathcal{F}^{\pi})_{\pi\geqslant\pi_0}$ is an fcga on the $\mathfrak{U}$-category $\mathcal{M}$, so is $(\mathcal{M}_\pi)_{\pi\geqslant\pi_0}$. Moreover, every useful (resp. regular) cardinal for the first filtration, is so for the second (called the} associated canonical \textit{fcg of the given} \texttt{[generalized]} \textit{cardinal filtration).}

The last assertion follows from the fact that if $\pi$ is useful for $(\mathcal{F}^{\pi})$, then we have $\mathcal{F}^{\pi}=\mathcal{M}_\pi$.

\textbf{4.7.16.} It seems to me that one can say that the canonical fcg associated with a given fcga
\msp{page 181}
is always ``better'' than the latter (if the latter is not canonical, i.e. if the two are not equal), because of the following two facts:
\begin{enumerate}[label=\alph*)]
\item precisely because it is \textit{canonical}, i.e. depends on no other choice than that of the cardinal $\pi_0$ alone, and
\item because there are a priori more useful cardinals and regular cardinals for the canonical fcg.
\end{enumerate}
So all in all, it seems to me that the \texttt{[generalized]} cardinal filtrations the only useful ones are those which are canonical.

\noindent\printedpage{62}% head: A. Grothendieck : Dérivateurs, Ch. XVIII, Édition des Universités de Montpellier II et Paris VII
% foot: Version 17 décembre 2025
Let us note that a cardinal filtration (in the sense of SGA 4 I 9, which I retain here) which is in addition \textit{accessible} (i.e. $\mathcal{F}^{\pi}\subset\mathcal{M}_\pi$ for every $\pi\geqslant\pi_0$), is nothing other than a canonical fcga, such that \textit{all} cardinals are regular (a fortiori useful) for the said filtration. I presume (without however having constructed an example) that the canonical \texttt{[generalized]} cardinal filtrations of categories of the form $\mathcal{M}=\Ind_{\pi_0}(C)$, $C$ a small \textit{Karoubian} category (so that $\pi_0$ is a useful cardinal) are not necessarily fc (not generalized, i.e. in the SGA sense), i.e. I presume that it happens that there are cardinals $\pi\geqslant\pi_0$ not regular, and even superfluous.

\noindent\msp{page 182}I must finally study what happens when we change the initial cardinal $\pi_0$ of an fcga.

\textbf{Proposition 4.7.17.} \textit{Let $(\mathcal{F}^{\pi})_{\pi\geqslant\pi_0}$ be an fcga on the $\mathfrak{U}$-category $\mathcal{M}$, and let $\pi'_0$ be a cardinal $\geqslant\pi_0$. Then $(\mathcal{F}^{\pi})_{\pi\geqslant\pi'_0}$ is an fcga on $\mathcal{M}$, which is canonical if the one we start with is canonical. For a cardinal $\pi\geqslant\pi'_0$, it is useful for $(\mathcal{F}^{\pi})_{\pi\geqslant\pi'_0}$ if and only if it is so for $(\mathcal{F}^{\pi})_{\pi\geqslant\pi_0}$. For it to be regular for $(\mathcal{F}^{\pi})_{\pi\geqslant\pi'_0}$, it} suffices \textit{that it be so for $(\mathcal{F}^{\pi})_{\pi\geqslant\pi_0}$, and this is also necessary when we assume in addition that $2^{\pi}\geqslant\pi'_1\overset{\text{def}}{=}\sup(\pi'_0,\cardred\mathcal{F}^{\pi'_0})$, a fortiori when $\pi\geqslant\pi'_1$.}

\textsc{Proof}. Conditions a) b) are trivially satisfied by the second filtration, being satisfied by the first, as is d) if it is by the first. The same is true of condition c) if $\pi=\pi^{\pi'_0}$, which indeed implies $\pi^{\pi_0}=\pi$, since in general
\[
\pi \;\leqslant\; \pi^{\pi_0} \;\leqslant\; \pi^{\pi'_0} \, .
\]
It remains to examine the cardinal bound
\[
\card I \;\longrightarrow\; \pi'^{\pi} \qquad \text{if } F\in\Ob\mathcal{F}^{\pi'} \ \underline{\text{and}}\ \pi'^{\pi}\geqslant\pi'_1 \, , % [sic] \longrightarrow -> \leqslant, a majoration (inequality) of cardinals is meant
\]
where this time $\pi'_1=\sup(\pi'_0,\cardred\mathcal{F}^{\pi'_0})$. In any event, we have $\pi'_1\geqslant\pi_1$, hence
\msp{page 183}
$\pi'^{\pi}\geqslant\pi'_1$ implies $\pi'^{\pi}\geqslant\pi_1$, hence we can apply condition c) for fcg for $(\mathcal{F}^{\pi})_{\pi\geqslant\pi_0}$, to conclude that it is still verified. So in short, \textit{all} the conditions for fcg (resp. fcga) are trivially verified for the \textit{restricted} fcg of the first. Likewise, the fact that a cardinal $\pi'\geqslant\pi'_0$ is useful is trivially independent of which of the two filtrations we consider, it depends only on the datum of the cardinal $\pi'$ and of the subcategory $\mathcal{F}^{\pi'}$ of $\mathcal{M}$. There remains the cardinal bound, in the case where we are interested in the \textit{regularity} of $\pi'$ for the fcga considered. The bound required is the same, but the precondition for the validity of this bound, namely
\[
\pi'^{\pi} \;\geqslant\; \pi'_1 \qquad \text{resp.} \qquad \pi'^{\pi} \;\geqslant\; \pi_1 \, ,
\]
is not the same. Since $\pi'_1\geqslant\pi_1$, we see that the precondition is less restrictive in the case of the first filtration, than of the second, whence the ``it suffices'' in the last assertion. But when $2^{\pi}\geqslant\pi'_1$, even the more restrictive one is automatic\msp{page 184}ally verified, since $\pi'^{\pi}\geqslant 2^{\pi}\geqslant\pi'_1$, and then the regularity of $\pi$ for the second filtration implies the regularity for the first, q.e.d. % [sic] le première -> la première, gender agreement (feminine)

\noindent\printedpage{63}% head: A. Grothendieck : Dérivateurs, Ch. XVIII, Édition des Universités de Montpellier II et Paris VII
% foot: Version 17 décembre 2025
\textbf{Corollary 4.7.18.} \textit{Suppose $\pi_0$ regular for the fcga (\footnote[77]{It suffices for this that $\pi_0$ be useful, i.e. $\mathcal{M}\xleftarrow{\ \sim\ }\Ind_{\pi_0}(\mathcal{F}^{\pi_0}=\mathcal{M}_{\pi_0})$, this implies regular, cf. 4.8.5.}). If we restrict the fcga $(\mathcal{F}^{\pi})_{\pi\geqslant\pi_0}$ to the cardinals $\pi\geqslant\pi'_0$, where $\pi'_0\geqslant\pi_1$, then we find an fcga where the cardinal bound (for $\pi=\pi^{\pi_0}$) is valid without the precondition $\pi'^{\pi}\geqslant\pi'_1$; more precisely, we then have}
\[
\pi'_1 \;\leqslant\; 2^{({\pi'_0}^{\pi_0})} \qquad \text{\textit{(a fortiori $\pi'_1\leqslant 2^{(2^{\pi'_0})}$)},} \tag*{(4.7.18.1)}
\]
\textit{hence the precondition $\pi'^{\pi}\geqslant\pi'_1$ is automatically verified if $\pi\geqslant\pi_0$, $\pi=\pi^{\pi'_0}$, $\pi'>\pi$ (cf. remark 4.7.7.2).}

This is therefore equivalent to the following cardinal bound (applied to the case of $\pi=\pi'_0$).

\textbf{Lemma 4.7.19.} \textit{Suppose $\pi_0$ regular for the fcga $(\mathcal{F}^{\pi})_{\pi\geqslant\pi_0}$. Then for every cardinal $\pi\geqslant\pi_1$ ($\overset{\text{def}}{=}\sup(\pi_0,\cardred\mathcal{F}^{\pi_0})$), we have}
\[
\cardred\mathcal{F}^{\pi} \;\leqslant\; 2^{(\pi^{\pi_0})} \, , \tag*{(4.7.19.1)}
\]
\textit{and if $\pi'\geqslant\pi_1$ and $F\in\Ob\mathcal{F}^{\pi'}$, we have}
\[
\cardred\mathcal{F}^{\pi}/F \;\leqslant\; \pi'^{(\pi^{\pi_0})} \qquad \text{(\footnotemark[78])} \, . \tag*{(4.7.19.2)}
\]
\footnotetext[78]{\textbf{N.B.} I do not know whether we can replace in these bounds $\pi^{\pi_0}$ by $\pi_0$.}

\noindent\msp{page 185}\textsc{Proof of the lemma}. Since $\pi_0$ is regular, putting $C=\mathcal{F}^{\pi_0}=\mathcal{M}_{\pi_0}$, we find $\mathcal{M}\simeq\Ind_{\pi_0}(C)$, hence we can assume $\mathcal{M}=\Ind_{\pi_0}(C)$, and apply the results proved for categories of this form. Since $\mathcal{F}^{\pi}\subset\mathcal{M}_\pi$, $\mathcal{F}^{\pi'}\subset\mathcal{M}_{\pi'}$, the bounds (4.7.9.1) and (4.7.9.2) are contained in % [sic] (4.7.9.1) et (4.7.9.2) -> (4.7.19.1) et (4.7.19.2), the majorations of Lemme 4.7.19 are meant
\[
\left\{
\begin{array}{rcll}
\cardred\mathcal{M}_\pi & \leqslant & 2^{(\pi^{\pi_0})} & \\
\cardred\mathcal{M}_\pi/F & \leqslant & \pi'^{(\pi^{\pi_0})} & \text{if } F\in\Ob\mathcal{M}_{\pi'},\ \text{with}\\
 & & & \pi'\geqslant\pi\geqslant\pi_1.
\end{array}
\right\}
\ \text{for } \pi\geqslant\pi_1 \qquad \text{(\footnotemark[79])}. \tag*{$(*)$}
\]
\footnotetext[79]{Compare with bounds $\cardred\Filt^{\pi}\leqslant 2^{\pi}$, $\cardred\Filt^{\pi}/F\leqslant\pi'^{\pi}$, if $F\in\Ob\Filt^{\pi'}$, with $\pi'\geqslant\pi\geqslant\pi_1$, bounds more pleasant than those of (4.7.19). We should all the same make sure on an example whether the nicer bounds might not be valid in the general case.}
The first of these bounds is nothing other than the second bound in 4.7.7.1. Let us note that we do indeed have
\[
2^{(\pi^{\pi_0})} \;\leqslant\; 2^{(2^{\pi})} \, ,
\]
which indeed follows from
\[
\pi^{\pi_0} \;\leqslant\; 2^{\pi} \, ,
\]
which we prove by noting that $\pi<2^{\pi}$, whence $\pi^{\pi_0}\leqslant(2^{\pi})^{\pi_0}=2^{\pi\cdot\pi_0}=2^{\pi}$ since $\pi\geqslant\pi_0$. It remains to establish the second bound
\[
\cardred\mathcal{M}_\pi/F \;\leqslant\; \pi'^{(\pi^{\pi_0})} \qquad \text{if } F\in\Ob\mathcal{M}_{\pi'},\ \text{with } \pi'\geqslant\pi\geqslant\pi_1 \, . \tag*{$(*)$}
\]
(We do not need it for 4.7.18, but it will probably be useful elsewhere.) We use $\mathcal{M}_{\pi'}\subset\Filt^{\pi'^{\pi_0}}$ if $\pi'\geqslant\pi_1$ (4.7.6.2), hence if $F\in\Ob\mathcal{M}_{\pi'}$, we have

\end{document}
